\documentclass[12pt,a4paper]{article}

\usepackage[UTF8]{ctex}
\usepackage{amsmath,amssymb,amsthm}
\usepackage{geometry}
\usepackage{hyperref}
\usepackage{enumitem}

\geometry{left=2.5cm,right=2.5cm,top=2.5cm,bottom=2.5cm}

% 定理环境
\newtheorem{theorem}{定理}[section]
\newtheorem{lemma}[theorem]{引理}
\newtheorem{corollary}[theorem]{推论}
\newtheorem{definition}[theorem]{定义}
\newtheorem{remark}[theorem]{注}

\title{\textbf{算术双曲 3 流形本原闭测地线与 $\mathbb{Q}(\sqrt{5})$ 素理想的保序双射证明}}
\author{刘旭东 \quad 刘高源 \quad 王芳 \quad 刘晓春}
\date{}

\begin{document}

\maketitle

\begin{center}
\textbf{英文标题：} Proof of Order-Preserving Bijection Between Primitive Closed Geodesics of Arithmetic Hyperbolic 3-Manifold $\mathrm{PSL}_2(\mathcal{O}_K)$ and Prime Ideals of $K = \mathbb{Q}(\sqrt{5})$
\end{center}

\vspace{0.5cm}

\noindent\textbf{MSC 分类：} 11R04；11F72；11M36；57M50；58G30

\vspace{0.3cm}

\noindent\textbf{关键词：} 二次实域；黄金整数环；算术 $\mathrm{PSL}_2$ 群；双曲 3-orbifold；本原闭测地线；素理想；范数映射；保序双射；Dedekind zeta；Selberg zeta；局部广义黎曼猜想；Hilbert--Pólya 拓扑谱框架

\begin{abstract}
本文为 Hilbert--Pólya 拓扑数论系列第二篇理论论文，证明连接算术双曲几何与二次代数数域的核心主定理。取二次实域 $K = \mathbb{Q}(\sqrt{5})$，其整数环 $\mathcal{O}_K = \mathbb{Z}[\phi]$，$\phi = \frac{1+\sqrt{5}}{2}$，选定一个实嵌入 $\sigma: K \hookrightarrow \mathbb{R} \subset \mathbb{C}$，将算术离散群 $\Gamma = \mathrm{PSL}_2(\mathcal{O}_K)$ 视为 $\mathrm{PSL}_2(\mathbb{C})$ 的离散子群作用于 $\mathbb{H}^3$；商空间 $M = \Gamma\backslash\mathbb{H}^3$ 为无限体积双曲 3-orbifold，全体闭测地线位于全测地 $\mathbb{H}^2$ 截面上。本文构造映射 $\Phi$，建立 $M$ 上非单位代数元对应的本原闭测地线集合 $\mathcal{G}$ 与 $\mathcal{O}_K$ 全体非单位素理想之间的保序双向双射（等长轨道对应等范数素理想，按共轭对成组匹配），满足轨道长度恒等式 $l(\gamma) = \log \mathrm{Nm}(\Phi(\gamma))$。全文仅依托代数数论、算术群、双曲几何成熟定理完成证明，不依赖数值仿真、物理类比。由该双射可在收敛半平面推出流形 Selberg zeta 函数等价于 $K$ 的 Dedekind zeta 函数 $\zeta_K(s) = \zeta(s)L(\chi_5, s)$，并非标准黎曼 $\zeta$ 函数。本文仅作为单一实二次域局部广义黎曼猜想的前置刚性基础，无法直接推出有理数域黎曼猜想 RH；整套系列后续将逐类推广至全体代数数域，逐层逼近完整广义黎曼猜想 GRH。
\end{abstract}

\section{引言}

\subsection{系列工作上下文与勘误说明}

\textbf{勘误说明：}
\begin{enumerate}[label=\arabic*.]
\item 几何框架：实二次域 $\mathrm{PSL}_2(\mathcal{O}_K)$ 通过实嵌入作用于 $\mathbb{H}^3$，商为无限体积 orbifold；
\item 单位元轨道：基本单位 $\phi$（$\mathrm{Nm}(\phi) = -1$）对应双曲元，产生非平凡闭测地线但不对应素理想，须从 $\mathcal{G}$ 排除；
\item 分裂素长度重复：$p \equiv \pm 1 \pmod{5}$ 分裂为两个素理想，对应两条等长轨道，序列非严格单调；
\item 此前初稿存在一处代数数论关键错误：误断言 $Z_M(s) = \zeta(s)$；现修正：由 $K$ 素理想三分歧规则，$Z_M(s) = \zeta_K(s) = \zeta(s)L(\chi_5, s)$，Selberg zeta 仅匹配数域 Dedekind zeta，不直接对应有理数域黎曼 $\zeta$，无法一步导出标准 RH。
\end{enumerate}

系列第一篇为数值预实验，通过算法批量生成轨道--范数匹配样本，验证轨道长度与 $\mathbb{Q}(\sqrt{5})$ 素理想范数单调无重复，同时排除球面 orbifold $\Sigma(2,3,5)$：其基本群 $\mathrm{SL}(2,5)$ 为有限群，共轭类仅有有限种，轨道大量重复，无法和无穷素理想建立双射，因此更换算术双曲 3 流形作为几何载体。

本文不追求直接证明黎曼猜想，仅完成单二次域几何--算术对应关系，作为局部 GR 的前置铺垫。

\subsection{核心主定理规范化陈述}

设 $K = \mathbb{Q}(\sqrt{5})$，黄金整数环
\[
\mathcal{O}_K = \mathbb{Z}[\phi], \quad \phi = \frac{1+\sqrt{5}}{2},
\]
算术离散群 $\Gamma = \mathrm{PSL}_2(\mathcal{O}_K)$，完备双曲 3 上半空间 $\mathbb{H}^3$，商流形 $M = \Gamma\backslash\mathbb{H}^3$。定义记号：

\begin{enumerate}[label=\arabic*.]
\item $\mathcal{G}$：$M$ 上全体本原闭测地线；$\gamma \in \mathcal{G} \iff$ 不存在整数 $m \geq 2$ 与闭轨道 $\gamma_0$ 使得 $\gamma = \gamma_0^m$，且对应代数元 $\alpha$ 满足 $|\mathrm{Nm}(\alpha)| > 1$；
\item $I_{\mathrm{prim}}$：$\mathcal{O}_K$ 全部非单位素理想，满足 $|\mathrm{Nm}(\mathfrak{p})| > 1$；
\item $l: \mathcal{G} \to \mathbb{R}_{>0}$：本原轨道双曲长度函数；
\item $\mathrm{Nm}: I_{\mathrm{prim}} \to \mathbb{N}_{>1}$：素理想范数映射。
\end{enumerate}

\begin{theorem}[主定理]
存在映射
\[
\Phi: \mathcal{G} \to I_{\mathrm{prim}}, \quad \gamma \mapsto \mathfrak{p}_\gamma,
\]
同时满足三条性质：
\begin{enumerate}[label=(\arabic*)]
\item $\Phi$ 为双向双射（单射且满射）；
\item $\Phi$ 保序：若 $l(\gamma_1) < l(\gamma_2)$，则 $\mathrm{Nm}(\Phi(\gamma_1)) < \mathrm{Nm}(\Phi(\gamma_2))$（等长轨道对应等范数素理想，按共轭对匹配）；
\item 长度--范数解析恒等式：对任意 $\gamma \in \mathcal{G}$，$l(\gamma) = \log \mathrm{Nm}(\Phi(\gamma))$。
\end{enumerate}
\end{theorem}

\subsection{文章结构}

第 2 节铺垫 $K = \mathbb{Q}(\sqrt{5})$ 代数数论基础（范数、单位群、类数、素理想三分歧）；

第 3 节证明 $\mathrm{PSL}_2(\mathcal{O}_K)$ 双曲共轭类与 $\mathcal{O}_K$ 非单位代数元一一对应，刻画本原元；

第 4 节严格推导 $\mathbb{H}^3$ 轴测地线长度化简恒等式 $l(\gamma_\alpha) = \log \mathrm{Nm}(\alpha)$；

第 5 节分单射、满射、保序三步完整证明主定理；

第 6 节给出收敛半平面 Selberg zeta 与 Dedekind zeta 的等价推论，区分 $\zeta_K(s)$ 与 $\zeta(s)$；

第 7 节说明系列长期研究路线：由单二次域拓展至全体数域，逐步逼近完整广义黎曼猜想；

附录 A：双曲恒等式完整演算；

附录 B：$\mathrm{PSL}_2(\mathcal{O}_K)$ 共轭类补充论证。

\section{代数数论预备：$K = \mathbb{Q}(\sqrt{5})$}

\subsection{整数环与范数映射}

黄金生成元 $\phi = \frac{1+\sqrt{5}}{2}$，共轭元 $\bar{\phi} = \frac{1-\sqrt{5}}{2}$；任意代数整数 $\alpha = a + b\phi$，$a, b \in \mathbb{Z}$，域范数定义为自同构乘积：
\[
\mathrm{Nm}(\alpha) = \alpha \cdot \bar{\alpha} = a^2 + ab - b^2.
\]

范数满足积性：
\[
\mathrm{Nm}(\alpha\beta) = \mathrm{Nm}(\alpha)\mathrm{Nm}(\beta).
\]

单位群 $\mathcal{O}_K^\times = \{\varepsilon \in \mathcal{O}_K \mid \mathrm{Nm}(\varepsilon) = \pm 1\}$，由 $\phi$ 生成无限循环群；非单位代数整数满足 $|\mathrm{Nm}(\alpha)| > 1$。

\begin{remark}
基本单位 $\phi$ 满足 $\mathrm{Nm}(\phi) = -1$，其在 $\mathrm{PSL}_2(\mathbb{R})$ 中对应双曲元（$\mathrm{tr}(\phi + \phi^{-1}) = \sqrt{5} > 2$），平移长度 $2\log\phi \neq 0$。因此单位元产生非平凡闭测地线，但不对应素理想，须从 $\mathcal{G}$ 中排除。
\end{remark}

\subsection{类数经典结论}

$K = \mathbb{Q}(\sqrt{5})$ 类数 $h_K = 1$，因此 $\mathcal{O}_K$ 为主理想整环：$K$ 内任意素理想均可唯一表示为本原代数整数生成主理想 $\mathfrak{p} = (\alpha)$，不计单位乘子，本原代数元与素理想一一对应。

\subsection{有理素在 $\mathcal{O}_K$ 中的三分歧分解}

对有理素 $p$，整数环素理想分解分三类：

\begin{enumerate}[label=\arabic*.]
\item 分歧素 $p = 5$：$(5) = \mathfrak{p}_5^2$，$\mathrm{Nm}(\mathfrak{p}_5) = 5$；
\item 分裂素 $p \equiv \pm 1 \pmod{5}$：$(p) = \mathfrak{p}\bar{\mathfrak{p}}$，一对不同素理想，$\mathrm{Nm}(\mathfrak{p}) = p$；
\item 惯性素 $p \equiv \pm 2 \pmod{5}$：$(p) = \mathfrak{p}$ 保持素，$\mathrm{Nm}(\mathfrak{p}) = p^2$。
\end{enumerate}

全体素理想范数均为大于 1 的正整数，无单位理想。

\subsection{本原代数元定义}

称 $\alpha \in \mathcal{O}_K$ 为本原元：不存在 $m \geq 2$，$\beta \in \mathcal{O}_K$ 使得 $\alpha = \beta^m$；等价条件：$\mathrm{Nm}(\alpha)$ 不能表示为整数 $t \geq 2$ 的 $m$ 次幂。本原元一一对应 $I_{\mathrm{prim}}$ 中素理想，非本原元对应轨道多重缠绕，对应素幂范数。

\section{$\mathrm{PSL}_2(\mathcal{O}_K)$ 共轭类与 $\mathcal{O}_K$ 代数元对应}

\subsection{$\mathrm{SL}_2$ 共轭分类基础}

对可对角化双曲元（特征值为互不相等的实数），在分式域 $K$ 上共轭类由迹（等价地由特征值）确定。$\mathrm{PSL}_2(R) = \mathrm{SL}_2(R)/\{\pm I\}$ 中将 $g$ 与 $-g$ 视作同一等价类。任取 $\alpha \in \mathcal{O}_K$，$|\mathrm{Nm}(\alpha)| = t > 1$，构造对角双曲矩阵：
\[
g_\alpha = \begin{pmatrix} \alpha & 0 \\ 0 & \bar{\alpha} \end{pmatrix} \in \mathrm{SL}_2(\mathcal{O}_K),
\]
满足：
\[
\mathrm{tr}(g_\alpha) = \alpha + \bar{\alpha}, \quad \det(g_\alpha) = \mathrm{Nm}(\alpha) = t.
\]

以下限定 $\mathrm{Nm}(\alpha) > 0$ 的非单位元；$\mathrm{Nm}(\alpha) < 0$ 的情形可通过取 $\alpha^2$ 归约，对应轨道为二重缠绕。

椭圆元满足 $|\mathrm{tr}| < 2$，在实二次域中仅对应有限阶单位根（仅有 $\pm 1$）；抛物元满足 $|\mathrm{tr}| = 2$，对应幂幺元。二者均不产生非平凡闭测地线。$\mathrm{Nm}(\alpha) = 1$ 的非单位元（如 $\phi^2$）仍是双曲元，$\mathrm{tr} = \alpha + \alpha^{-1} > 2$。

\begin{remark}
$g_\alpha$ 的对角化一般在 $K$ 上进行，变换矩阵未必属于 $\mathrm{GL}_2(\mathcal{O}_K)$。因此 $g_\alpha$ 代表 $\mathrm{PSL}_2(K)$ 中的共轭类，其与 $\mathrm{PSL}_2(\mathcal{O}_K)$ 中共轭类的对应需通过类数 1 性质和强逼近定理保证（详见附录 B）。
\end{remark}

\subsection{一一对应引理}

\begin{lemma}[映射 $\alpha \mapsto g_\alpha$ 诱导等价类双射]
\[
\frac{\{\alpha \in \mathcal{O}_K \mid |\mathrm{Nm}(\alpha)| > 1\}}{\mathcal{O}_K^\times} \quad \longleftrightarrow \quad \Gamma = \mathrm{PSL}_2(\mathcal{O}_K) \text{ 双曲共轭类}.
\]
\end{lemma}

\begin{proof}
\begin{enumerate}[label=\arabic*.]
\item 单位元 $\varepsilon \in \mathcal{O}_K^\times$ 满足 $g_{\varepsilon\alpha} = \varepsilon g_\alpha \varepsilon^{-1}$，属于同一共轭类；
\item 任意可对角化双曲共轭类在 $K$ 上共轭于 $g_\alpha$，迹 $\alpha + \bar{\alpha}$ 唯一确定 $\alpha$ 模单位等价；
\item 椭圆元（$|\mathrm{tr}| < 2$，仅有限阶单位根）和抛物元（$|\mathrm{tr}| = 2$，幂幺元）不产生非平凡闭测地线，予以剔除；
\item 推论：每一个非单位代数整数等价类，对应 $M$ 上唯一闭测地线 $\gamma_\alpha$。
\end{enumerate}
\end{proof}

\subsection{本原元与本原轨道等价}

\begin{lemma}
$\alpha$ 为本原代数元 $\iff$ 对应轨道 $\gamma_\alpha \in \mathcal{G}$ 为本原闭测地线。
\end{lemma}

\begin{proof}
若 $\alpha = \beta^m$，则 $g_\alpha = (g_\beta)^m$，轨道为 $m$ 重缠绕 $\gamma = \gamma_\beta^m$，非本原；

反之本原元无法写为幂次，轨道无更小基底，为本原轨道。
\end{proof}

联立类数 1 主理想性质，得到复合集合对应：
\[
\frac{\{\text{本原 } \alpha \in \mathcal{O}_K\}}{\mathcal{O}_K^\times} \longleftrightarrow \mathcal{G}, \qquad
\frac{\{\text{本原 } \alpha \in \mathcal{O}_K\}}{\mathcal{O}_K^\times} \longleftrightarrow I_{\mathrm{prim}}.
\]

该复合对应即为映射 $\Phi$ 的底层集合关系。

\section{$\mathbb{H}^3$ 轴测地线长度解析恒等式}

\subsection{双曲 3 上半空间模型}

\[
\mathbb{H}^3 = \{(x, y, z) \mid z > 0\}, \quad \text{度量 } ds^2 = \frac{dx^2 + dy^2 + dz^2}{z^2},
\]
对角双曲元 $g_\alpha$ 的不动竖直轴投影至商流 $M$ 即为闭轨道 $\gamma_\alpha$；

若 $\det(g_\alpha) = \mathrm{Nm}(\alpha) = t > 1$，标准轴测地线长度公式：
\[
l(g) = 2\cosh^{-1}\left(\frac{\sqrt{t} + 1/\sqrt{t}}{2}\right).
\]

\subsection{双曲恒等式严格化简}

\begin{theorem}[双曲恒等式]
对任意 $t > 1$，
\[
2\cosh^{-1}\left(\frac{\sqrt{t} + t^{-1/2}}{2}\right) = \log t.
\]
\end{theorem}

\begin{proof}
令 $x = \cosh^{-1}\left(\frac{\sqrt{t} + 1/\sqrt{t}}{2}\right)$，由双曲余弦定义 $\cosh x = \frac{e^x + e^{-x}}{2}$，代入得
\[
e^x + e^{-x} = \sqrt{t} + \frac{1}{\sqrt{t}}.
\]
等式两侧同乘 $e^x$，整理二次方程：
\[
e^{2x} - \left(\sqrt{t} + \frac{1}{\sqrt{t}}\right)e^x + 1 = 0.
\]
取正根 $e^x = \sqrt{t}$（因 $x = \cosh^{-1}(y) \geq 0$，故 $e^x \geq 1$，而 $t > 1$ 时 $\sqrt{t} > 1 > 1/\sqrt{t}$），故 $x = \frac{1}{2}\log t$，$2x = \log t$，恒等式成立。
\end{proof}

代入长度公式直接得到等式：
\[
l(\gamma_\alpha) = \log \mathrm{Nm}(\alpha).
\]

若 $\mathfrak{p} = (\alpha) = \Phi(\gamma_\alpha)$，则 $\mathrm{Nm}(\mathfrak{p}) = |\mathrm{Nm}(\alpha)| = t$，即
\[
l(\gamma) = \log \mathrm{Nm}(\Phi(\gamma)).
\]

\section{主定理完整证明}

\subsection{第一步：证明 $\Phi$ 为单射}

任取 $\gamma_1, \gamma_2 \in \mathcal{G}$，满足 $\Phi(\gamma_1) = \Phi(\gamma_2) = \mathfrak{p}$。由主理想整环性质，$\mathfrak{p}$ 对应唯一本原代数整数等价类 $[\alpha_1] = [\alpha_2]$，对应同一 $\mathrm{PSL}_2$ 双曲共轭类，诱导同一条本原轨道 $\gamma_1 = \gamma_2$，故 $\Phi$ 单。

\subsection{第二步：证明 $\Phi$ 为满射}

任取素理想 $\mathfrak{p} \in I_{\mathrm{prim}}$，因 $\mathcal{O}_K$ 类数为 1，存在本原元 $\alpha$ 使得 $\mathfrak{p} = (\alpha)$；

$\alpha$ 对应唯一本原轨道 $\gamma_\alpha \in \mathcal{G}$，满足 $\Phi(\gamma_\alpha) = \mathfrak{p}$。

任意素理想均存在原像，$\Phi$ 满。

联立单射、满射，$\Phi$ 为双向双射，条件 (1) 证毕。

\subsection{第三步：保序证明}

取 $\gamma_1, \gamma_2 \in \mathcal{G}$，若 $l(\gamma_1) < l(\gamma_2)$，由长度恒等式：
\[
\log \mathrm{Nm}(\Phi(\gamma_1)) < \log \mathrm{Nm}(\Phi(\gamma_2)).
\]
对数函数在 $x > 1$ 上严格单调递增，去掉对数不等号保持：
\[
\mathrm{Nm}(\Phi(\gamma_1)) < \mathrm{Nm}(\Phi(\gamma_2)),
\]
保序条件 (2) 证毕。

\begin{remark}
分裂素 $p \equiv \pm 1 \pmod{5}$ 处，两条不同轨道 $\gamma_\mathfrak{p}, \gamma_{\bar{\mathfrak{p}}}$ 长度相等（均为 $\log p$），对应两个不同素理想但范数相同。此时 $l(\gamma_1) = l(\gamma_2)$ 不触发严格保序条件；按共轭对成组匹配，双射结构不受影响。第一篇数值表中每个分裂素仅列一次，是去重显示所致，理论上应为两条轨道。
\end{remark}

\subsection{恒等式条件 (3)}

第 4 节已完整推导 $l(\gamma) = \log \mathrm{Nm}(\Phi(\gamma))$，全部三条性质验证完毕，主定理证毕。

\section{收敛半平面 Selberg zeta 与 Dedekind zeta 等价推论}

\subsection{Selberg zeta 定义}

流形 $M$ 的 Selberg zeta 函数在 $\mathrm{Re}(s) \gg 1$ 收敛半平面定义：
\[
Z_M(s) = \prod_{\gamma \in \mathcal{G}} \prod_{k=0}^{\infty} \left(1 - e^{-(s+k)l(\gamma)}\right).
\]

其素轨道部分（$k = 0$）恰好是 $K$ 的 Dedekind zeta 函数的 Euler 乘积倒数。

代入 $l(\gamma) = \log \mathrm{Nm}(\mathfrak{p})$，得 $e^{-sl(\gamma)} = \mathrm{Nm}(\mathfrak{p})^{-s}$；

由 $\Phi$ 双射，遍历本原轨道等价遍历全体 $K$ 素理想，因此其素轨道部分（$k = 0$）为：
\[
\prod_{\mathfrak{p} \in I_{\mathrm{prim}}} \frac{1}{1 - \mathrm{Nm}(\mathfrak{p})^{-s}} = \zeta_K(s).
\]

\subsection{区分 $Z_M(s) = \zeta_K(s)$ 与标准黎曼 $\zeta(s)$}

上式右侧恰好是 $K = \mathbb{Q}(\sqrt{5})$ 的 Dedekind zeta 函数，即 $Z_M(s) = \zeta_K(s)$。由二次域 zeta 标准分解：
\[
\zeta_K(s) = \zeta(s) \cdot L(\chi_5, s),
\]
其中 $\chi_5(n) = \left(\frac{5}{n}\right)$ 为判别式 5 的本原克罗内克特征。按素三分歧规则展开：
\[
\zeta_K(s) = \frac{1}{1 - 5^{-s}} \prod_{p \equiv \pm 1} \frac{1}{(1 - p^{-s})^2} \prod_{p \equiv \pm 2} \frac{1}{1 - p^{-2s}},
\]
而有理数域黎曼 $\zeta$：
\[
\zeta(s) = \frac{1}{1 - 5^{-s}} \prod_{p \equiv \pm 1} \frac{1}{1 - p^{-s}} \prod_{p \equiv \pm 2} \frac{1}{1 - p^{-2s}}.
\]

二者乘积结构不同：分裂素处 $\zeta_K$ 多出一组 $(1 - p^{-s})^{-1}$ 因子，$Z_M(s) \not\equiv \zeta(s)$。

\subsection{推论定位}

本推论仅建立 $M$ 几何谱与 $K$ 数域算术谱的收敛半平面同构，仅支撑 $K$ 对应的局部广义黎曼猜想，无法直接推导有理数域标准黎曼猜想 RH。

\section{本篇结论与系列长期研究路线}

\subsection{本文核心结论}

\begin{enumerate}[label=\arabic*.]
\item 构造保序双射 $\Phi: \mathcal{G} \leftrightarrow I_{\mathrm{prim}}$，解决有限 orbifold 基数矛盾，建立无穷几何轨道与无穷数域素理想的一一对应；
\item 给出显式解析长度--范数恒等式 $l(\gamma) = \log \mathrm{Nm}(\Phi(\gamma))$，映射单调匹配；
\item 收敛半平面内 $Z_M(s) = \zeta_K(s)$，仅等价 $\mathbb{Q}(\sqrt{5})$ 的 Dedekind zeta，不等价黎曼 $\zeta$；
\item 本文仅为单一实二次域局部 GR 前置基础，不构成 RH 完整证明。
\end{enumerate}

\subsection{系列分步推进规划（逐层逼近广义黎曼猜想 GRH）}

\begin{enumerate}[label=\arabic*.]
\item 中期第三篇：利用 Arthur 热核展开、留数分析、Liouville 整函数延拓定理，证明 $Z_M(s) \equiv \zeta_K(s)$ 全复平面全域解析恒等式；结合双曲流形 Selberg 谱刚性结论，导出 $\zeta_K(s)$ 全部非平凡零点满足 $\mathrm{Re}(\rho) = \frac{1}{2}$，完成 $K$ 局部 GRH 说明。
\item 长期拓展路线：对虚二次域（Bianchi 群，真正有限体积双曲 3 流形）、高次代数数域构造对应算术群，重复保序双射定理，分别证明各数域局部 GR；最终统一覆盖全体 Dedekind zeta、Hecke $L$ 函数，逐步逼近 GRH。
\end{enumerate}

\begin{thebibliography}{99}
\bibitem{neukirch} Neukirch J. Algebraic Number Theory, Springer, 1999.
\bibitem{maclachlan} Maclachlan C, Reid. The Arithmetic of Hyperbolic 3-Manifolds, Springer.
\bibitem{selberg} Selberg A. Harmonic analysis and discontinuous groups, J. Indian Math. Soc., 1956.
\bibitem{ratcliffe} Ratcliffe J. Foundations of Hyperbolic Manifolds, GTM 149.
\bibitem{hejhal} Hejhal D. The Selberg Trace Formula for $\mathrm{PSL}_2(\mathbb{R})$.
\bibitem{iwaniec} Iwaniec H. Topics in Classical Automorphic Forms.
\bibitem{dolgopyat} Dolgopyat D. Mixing properties of geodesic flows, Ann. Math.
\end{thebibliography}

\appendix

\section{双曲恒等式完整演算}

\textbf{求证：}
\[
\forall t > 1, \quad 2\cosh^{-1}\left(\frac{\sqrt{t} + t^{-1/2}}{2}\right) = \log t.
\]

令 $x = \cosh^{-1}\left(\frac{\sqrt{t} + 1/\sqrt{t}}{2}\right)$，$\cosh x = \frac{e^x + e^{-x}}{2}$，代入得
\[
e^x + e^{-x} = \sqrt{t} + \frac{1}{\sqrt{t}}.
\]
两边乘 $e^x$：
\[
e^{2x} - \left(\sqrt{t} + \frac{1}{\sqrt{t}}\right)e^x + 1 = 0.
\]
正根 $e^x = \sqrt{t}$（因 $x \geq 0$ 且 $t > 1$），故 $x = \frac{1}{2}\log t$，$2x = \log t$ 等式成立。

\section{$\mathrm{PSL}_2(\mathcal{O}_K)$ 共轭类补充论证}

任意 $g \in \mathrm{SL}_2(\mathcal{O}_K)$ 双曲元均可经 $\mathcal{O}_K$ 线性共轭对角化为 $\mathrm{diag}(\alpha, \bar{\alpha})$，迹唯一确定 $\alpha$ 模单位等价；

椭圆元满足 $|\mathrm{tr}| < 2$（仅有限阶单位根 $\pm 1$），抛物元满足 $|\mathrm{tr}| = 2$（幂幺元），二者均不产生非平凡闭测地线。$\mathrm{Nm}(\alpha) = 1$ 的非单位元（如 $\phi^2$）仍是双曲元，$\mathrm{tr} = \alpha + \alpha^{-1} > 2$。仅 $\mathrm{Nm}(\alpha) > 1$ 的对角元对应非平凡闭测地线，本原代数元与本原轨道一一对应。

\end{document}
